\loai{Tính công cơ học}
\begin{vidu} %[3P3-13.1B][Loai1] 
	Một động cơ điện xoay chiều tiêu thụ công suất $1\ kW$ và có hiệu suất 80\%. Công cơ học mà động cơ sinh ra trong 30 phút bằng
	\choice
	{\True $1440\ kJ$}
	{$1440\ kWh$}
	{$2250 \ kJ$}
	{$1440\ kW$}
	\loigiai{
%		\Binhluan{
		\begin{itemize}
%			\item Động cơ điện chuyển hóa điện năng (năng lượng toàn phần) $\Rightarrow $ cơ năng (năng lượng có ích) $+$ nhiệt năng (năng lượng hao phí).
			\item Công suất điện tiêu thụ chính là công suất toàn phần của động cơ $\Rightarrow $ công suất cơ học là công suất có ích
			$P_{ci}=P_{tp}.H = 800\ W$.
			\item Công cơ học trong thời gian 30 phút là $A =P_{ci}.t = 1440000\ J = 1440\ kJ$.
	\end{itemize}
%}
%{Học sinh phân biệt rõ các loại công suất ứng với động cơ điện}
}
\end{vidu}
\begin{vidu} %[3P3-13.1B][Loai1] 
	Một động cơ điện xoay chiều tiêu thụ công suất 1,5 kW và có hiệu suất 80\%. Nhiệt lượng do động cơ toả ra trong 30 phút khi động cơ hoạt động là
	\choice
	{$2,16.10^3$ J}
	{$4,32.10^3$ J}
	{\True $5,4.10^5$ J}
	{$2,16.10^6$ J}
	\loigiai{
%\Binhluan{
		Nhiệt lượng do động cơ toả ra trong 30 phút:  $A_{\text{toả nhiệt}}=0,2.P.t=0,2.1500.1800=5,4.10^5\ J$.
%}
%{Đây chính là phần năng lượng hao phí}
	}
\end{vidu}
\loai{Tỉ số}
\begin{vidu} %[3P3-13.1K][Loai2] 
	Một động cơ điện tiêu thụ công suất điện 110 W, sinh ra công suất cơ học bằng 88 W. Tỉ số của công suất cơ học với công suất hao phí ở động cơ bằng
	\choice
	{3}
	{\True 4}
	{2}
	{5}
	\loigiai{
%\Binhluan{
		Ta có: $P={{P}_{C}}+{{P}_{hp}}\Rightarrow {{P}_{hp}}=P-{{P}_{C}}=22\ W\Rightarrow \dfrac{{{P}_{C}}}{{{P}_{hp}}}=4$.
%}
%{
%Khi ta nói đến hiệu suất thì có thể tính theo năng lượng hoặc công suất đều như nhau.
%}
	}
\end{vidu}
\loai{Hiệu suất}
\begin{vidu} %[3P3-13.1K][Loai3] 
	\nguon{ĐH - 2012} Một động cơ điện xoay chiều hoạt động bình thường với điện áp hiệu dụng 220\ V, cường độ dòng điện hiệu dụng 0,5 A và hệ số công suất của động cơ là 0,8. Biết rằng công suất hao phí của động cơ là 11 W. Hiệu suất của động cơ (tỉ số giữa công suất hữu ích và công suất tiêu thụ toàn phần) là
	\choice
	{80\%}
	{90\%}
	{92,5\%}
	{\True 87,5\%}
	\loigiai{
%\Binhluan{
		\begin{itemize}
			\item Công suất tiêu thụ toàn phần:  $P=U.I.\cos\varphi =88\ W$.
			\item Công suất hữu ích: $P_{hi} = P - P_{hp} = 88-11 = 77\ W$.
			\item Hiệu suất của động cơ: $H = \dfrac{{P}_{ci}}{{P}_{hp}}=\dfrac{77}{88}=87,5\%$.
		\end{itemize}
%}
%{Công suất tiêu thụ điện: $P=UI\cos \varphi$ chính là công suất toàn phần}
	}
\end{vidu}
\begin{vidu} %[3P3-13.1B][Loai3] 
	Một động cơ điện xoay chiều hoạt động bình thường với điện áp hiệu dụng bằng $200\ V$ và cường độ dòng điện hiệu dụng bằng 0,5 A. Nếu công suất tỏa nhiệt trên dây quấn là 8 W và hệ số công suất của động cơ là 0,8 thì hiệu suất của động cơ là
	\choice
	{$86\%$}
	{\True $90\%$}
	{$75\%$}
	{$80\%$}
	\loigiai{
	%\Binhluan{
		\begin{itemize}
			\item Công suất tiêu thụ của động cơ:  $P=UI\cos\varphi  = 80\ W$.
			\item Công suất tỏa nhiệt $\Delta P$ trên dây quấn là công suất hao phí nên hiệu suất của động cơ
			\begin{align*}
				H=\dfrac{P-\Delta P}{P}=1-\dfrac{8}{80}=0,9=90\%.
			\end{align*}
		\end{itemize}
%}
%{Công suất toả nhiệt chính là công suất hao phí}
	}
\end{vidu}
\loai{Tính điện áp hai đầu động cơ}
\begin{vidu} %[3P3-13.1K][Loai4] 
	Một động cơ điện xoay chiều sản ra một công suất cơ học 8,5 kW và có hiệu suất 88\%. Xác định điện áp hiệu dụng ở hai đầu động cơ biết dòng điện có giá trị hiệu dụng 50 A và trễ pha so với điện áp hai đầu động cơ là $\dfrac{\pi}{12}$.
	\choice
	{331 V}
	{\True 200 V}
	{231 V}
	{565 V}
		\loigiai{
	%\Binhluan{
			Ta có: $P = \dfrac{{{P_i}}}{H} = UI\cos \varphi \Rightarrow U = \dfrac{{{P_i}}}{{HI\cos \varphi }} = \dfrac{{8,{{5.10}^3}}}{{0,88.50.\cos \dfrac{\pi }{{12}}}} = 200\ V.$ 
%	}
%{
%Ta không có công thức tính trực tiếp công suất cơ học (công suất có ích) thông qua dòng điện mà chỉ có thể tính toán thông qua hiệu suất của động cơ.
%}
		}
\end{vidu}
\loai{Tính cường độ dòng điện qua động cơ}
\begin{vidu} %[3P3-13.2K][Loai5] 
	Một động cơ điện xoay chiều có điện trở dây cuốn là 32 $\Omega$, mạch điện có điện áp hiệu dụng 200 V thì sản ra công suất cơ học 43 W. Biết hệ số công suất của động cơ là 0,9 và công suất hao phí nhỏ hơn công suất cơ học. Cường độ dòng hiệu dụng chạy qua động cơ là
	\choice
	{\True 0,25 A}
	{5,375 A}
	{0,225 A}
	{17,3 A}
	\loigiai{
%\Binhluan{
		\begin{itemize}
			\item \noindent $UI\cos \varphi = {P_i} + {I^2}r \Rightarrow 200.I.0,9 = 43 + {I^2}.32 \Rightarrow \left[ \begin{array}{l}
			I_1 = 5,375\ A\\
			I_2 = 0,25\ A
			\end{array} \right.$
			\item \noindent Ta chọn nghiệm $I_2 = 0,25$ A vì với nghiệm thứ nhất cho công suất hao phí lớn hơn công suất có ích.
		\end{itemize} 
%}
%{
%Trong động cơ điện một phần năng lượng điện đã biến thành năng lượng từ bên trong cuộn dây và tiếp tục biến thành năng lượng cơ học.
%}
	}
\end{vidu}



